% Heteroclinic-cycle figure, Figure 2.1(b) of the paper.
%
% Reads the tables ../results/hc-*.dat written by heteroclinic_cycle.py.
% Build from the repository root with
%
%     python3 heteroclinic_cycle.py --replot
%
% which runs latexmk in this directory and moves the PDF to ../figures/.
\documentclass[border=2pt]{standalone}
\usepackage{amsmath,amssymb}
\usepackage{pgfplots}
\pgfplotsset{compat=1.18}
\usetikzlibrary{arrows.meta,decorations.markings}

\definecolor{cycle}{HTML}{1F4E79}
\definecolor{unstable}{HTML}{C1121F}
\definecolor{stable}{HTML}{2A9D8F}
\definecolor{orbit}{HTML}{888888}

\tikzset{
    flow/.style={decoration={markings, mark=at position #1 with {\arrow{Stealth[length=5pt]}}},
        postaction=decorate},
    flow/.default=0.5,
}

\begin{document}
\begin{tikzpicture}
\begin{axis}[
        view={135}{32}, hide axis, width=7.5cm, height=7.5cm, scale only axis,
        xmin=-0.5, xmax=1.35, ymin=-0.5, ymax=1.35, zmin=-0.5, zmax=1.35,
        clip=false,
    ]

    % W^s(\bar x_2) cap X: the face x_3 = 0, with a few of its orbits
    \def\L{1.25}
    \addplot3[fill=stable, fill opacity=0.12, draw=none] coordinates {(0,0,0) (\L,0,0) (\L,\L,0) (0,\L,0)};
    \pgfplotsinvokeforeach{1,2,3,4}{
        \addplot3[stable, line width=0.5pt, no marks, flow=0.4] table[x=x1, y=x2, z=x3] {../results/hc-stable-#1.dat};
    }

    % coordinate axes
    \draw[-{Stealth[length=4pt]}, gray] (axis cs:0,0,0) -- (axis cs:1.6,0,0) node[pos=1, below left=-1pt] {$x_1$};
    \draw[-{Stealth[length=4pt]}, gray] (axis cs:0,0,0) -- (axis cs:0,1.5,0) node[pos=1, right] {$x_2$};
    \draw[-{Stealth[length=4pt]}, gray] (axis cs:0,0,0) -- (axis cs:0,0,1.4) node[pos=1, above] {$x_3$};

    % interior orbit
    \addplot3[orbit, line width=0.5pt, no marks, flow=0.93] table[x=x1, y=x2, z=x3] {../results/hc-orbit.dat};

    % the cycle, with Gamma_1^+ = W_+^u(\bar x_1) highlighted
    \addplot3[unstable, line width=1.1pt, no marks, flow] table[x=x1, y=x2, z=x3] {../results/hc-gamma-1.dat};
    \pgfplotsinvokeforeach{2,3}{
        \addplot3[cycle, line width=1.1pt, no marks, flow] table[x=x1, y=x2, z=x3] {../results/hc-gamma-#1.dat};
    }

    % the other branch of W^u(\bar x_1), which leaves X
    \addplot3[unstable, dashed, line width=0.8pt, no marks, flow=0.6] table[x=x1, y=x2, z=x3] {../results/hc-minus.dat};

    % saddles
    \addplot3[only marks, mark=*, mark size=2pt, mark options={fill=white, draw=black}]
        coordinates {(1,0,0) (0,1,0) (0,0,1)};
    \node[above left]  at (axis cs:1,0,0) {$\bar x_1$};
    \node[above right] at (axis cs:0,1,0) {$\bar x_2$};
    \node[left]        at (axis cs:0,0,1) {$\bar x_3$};

    % labels of the connections
    \node[unstable, below=2pt, fill=stable!12, inner sep=1pt] at (axis cs:0.45,0.55,0) {$\Gamma_1^+$};
    \node[unstable, above] at (axis cs:1.58,-0.52,0) {$W^u(\bar x_1)\setminus\mathcal X$};
    \node[stable, below] at (axis cs:1.25,1.25,0) {$W^s(\bar x_2)$};
    \node[cycle] at (axis cs:0,0.75,0.75) {$\Gamma_2^+$};
    \node[cycle] at (axis cs:0.75,0,0.75) {$\Gamma_3^+$};
\end{axis}
\end{tikzpicture}
\end{document}
